\subsection{FILTRATION OF STANDARD GROUPS}

We again use the notation of Definition 1. We choose a number \(a > 1\) and a real valuation \(v\) of \(\mathbf{K}\) such that \(|x| = a^{-v(x)}\) for all \(x \in \mathbf{K}\) (Commutative Algebra, Chapter VI, §6, Proposition 3). If \(a\) is a non-zero (and hence open) ideal of A contained in m, let \(\mathrm{G}(a)\) denote the set of elements of G whose coordinates belong to \(a\) . If \(\lambda \in \mathbf{R}\) , let \(a_{\lambda}\) (resp. \(a_{\lambda}^{+}\) ) denote the set of \(x \in \mathbf{K}\) such that \(v(x) \geqslant \lambda\) (resp. \(v(x) > \lambda\) ); then \(a_0 = A\) , \(a_0^+ = m\) . For \(x = (x_1, \ldots, x_r) \in G\) , we shall write

\[
\omega (x) = \inf (v (x _ {1}), \dots , v (x _ {r})).\tag{5}
\]

PROPOSITION 5. Let G be a standard group.

\begin{enumerate}
\def\labelenumi{(\roman{enumi})}
\item
  If \(\mathfrak{a}\) is a non-zero ideal of \(\mathbf{A}\) contained in \(\mathfrak{m}\) , \(\mathbf{G}(\mathfrak{a})\) is an open normal subgroup of \(\mathbf{G}\) .
\item
  The \(\mathbf{G}(\mathfrak{a}_{\lambda})\) , for \(\lambda > 0\) , form a fundamental system of neighbourhoods of \(e\) in \(\mathbf{G}\) .
\item
  Suppose that \(\mathfrak{a}_{\lambda} \subset \mathfrak{a}\) for \(\lambda \geqslant \lambda_0\) and let the \(G(\mathfrak{a}) / G(\mathfrak{a}_{\lambda})\) , for \(\lambda \geqslant \lambda_0\) , be given the discrete topology. Then the topological group \(G(\mathfrak{a})\) is the inverse limit of the groups \(G(\mathfrak{a}) / G(\mathfrak{a}_{\lambda})\) .
\item
  Let \(\mathfrak{a},\mathfrak{b}\) be non-zero ideals of A contained in \(\mathfrak{m}\) such that \(\mathfrak{a} \supset \mathfrak{b} \supset \mathfrak{a}^2\) . The mapping \(x \mapsto (x_1 \bmod \mathfrak{b}, \ldots, x_r \bmod \mathfrak{b})\) of G(a) into (a/b) × · · · × (a/b) defines on passing to the quotient an isomorphism of the group G(a)/G(b) onto the additive group (a/b) × · · · × (a/b).
\end{enumerate}

If \(x \in G\) and \(y \in G(a)\) , the coordinates of \(x\) and \(x.y\) are equal modulo \(a\) . Hence, for \(x'\) , \(x''\) in \(G\) and \(y'\) , \(y''\) in \(G(a)\) , the coordinates of \(x'.x''\) and \((x'.y')\) . \((x''.y'')\) are equal modulo \(a\) . This proves (i).

\begin{enumerate}
\def\labelenumi{(\roman{enumi})}
\setcounter{enumi}{1}
\item
  is obvious.
\item
  follows from the above and General Topology, Chapter III, § 7, Proposition 2.
\end{enumerate}

If \(x \in G(a)\) and \(y \in G(a)\) , the coordinates of \(x, y\) are congruent to those of \(x + y\) modulo \(G(a^2)\) by formula (4) of § 5. This proves (iv).

COROLLARY. Suppose that K is locally compact and let \(q = \operatorname{Card}(\mathbf{A} / \mathfrak{m})\) .

\begin{enumerate}
\def\labelenumi{(\roman{enumi})}
\item
  If \(a = m^{a}\) and \(b = m^{b}\) with \(b \geqslant a \geqslant 1\) , \(\mathrm{G}(a)/\mathrm{G}(b)\) is a p-group of cardinal \(q^{r(b-a)}\) .
\item
  G(a) is an inverse limit of p-groups.
\end{enumerate}

The number of elements in \(\mathrm{G}(\mathfrak{a})/\mathrm{G}(\mathfrak{b})\) is \((\mathrm{Card}(\mathfrak{a}/\mathfrak{b}))^{r}\) ; if \(b=a+1\) , a/b is a 1-dimensional vector space over A/m, whence (i) in this case; the general case follows by induction on b-a. Assertion (ii) follows from (i) and Proposition 5 (iii).

PROPOSITION 6. Let \(\mathfrak{a}, \mathfrak{b}, \mathfrak{c}, \mathfrak{c}'\) be non-zero ideals of \(\mathbf{A}\) contained in \(\mathfrak{m}\) such that

\[
\mathfrak {c} ^ {\prime} \subset \mathfrak {c}, \quad \mathfrak {a b} \subset \mathfrak {c}, \quad \mathfrak {a b} ^ {2} \subset \mathfrak {c} ^ {\prime}, \quad \mathfrak {a} ^ {2} \mathfrak {b} \subset \mathfrak {c} ^ {\prime}.
\]

If \(x \in \mathrm{G}(\mathfrak{a})\) and \(y \in \mathrm{G}(\mathfrak{b})\) , then \(x^{[-1]}y^{[-1]}x.y, x.y.x^{[-1]}y^{[-1]}\) and \([x,y]\) belong to \(\mathrm{G}(\mathfrak{c})\) and are congruent modulo \(\mathrm{G}(\mathfrak{c}')\) .

By § 5, no. 2, Proposition 1, there exist \(c_{\alpha \beta} \in \mathbf{A}^r\) such that

\[
x ^ {[ - 1 ]}. y ^ {[ - 1 ]}. x. y - [ x, y ] = \sum_ {| \alpha | + | \beta | \geqslant 3} c _ {\alpha \beta} x ^ {\alpha} y ^ {\beta}.
\]

If \(x = 0\) or \(y = 0\) , then \(x^{[-1]}y^{[-1]}x.y - [x,y] = 0\) ; hence \(c_{0\beta} = c_{\alpha 0} = 0\) . On the other hand, the conditions

\[
x \in \mathrm{G} (\mathfrak {a}), \quad y \in \mathrm{G} (\mathfrak {b}), \quad | \alpha | \geqslant 1, \quad | \beta | \geqslant 1, \quad | \alpha | + | \beta | \geqslant 3
\]

imply

\[
c _ {\alpha \beta} x ^ {\alpha} y ^ {\beta} \in \mathrm{G} (a ^ {2} b + a b ^ {2}) \subset \mathrm{G} (c ^ {\prime})
\]

and hence \(x^{[-1]} \cdot y^{[-1]}x \cdot y - [x, y] \in \mathrm{G}(\mathfrak{c}')\) . We see similarly that

\[
x. y. x ^ {[ - 1 ]}. y ^ {[ - 1 ]} - [ x, y ] \in \mathrm{G} (\mathfrak {c} ^ {\prime}).
\]

Finally, by § 5, formula (13), \([x,y] \in G(\mathfrak{ab}) \subset G(\mathfrak{c})\) .

PROPOSITION 7. (i) The family \((\mathrm{G}(\mathfrak{a}_{\lambda}))\) is a central filtration on \(\mathbf{G}\) (Chapter II, § 4, no. 4, Definition 2).

\[
\text { (ii) } For \lambda \in \mathbf {R} _ {+} ^ {*}, G (a _ {\lambda}) = \{x \in G | \omega (x) \geqslant \lambda \}, G (a _ {\lambda} ^ {+}) = \{x \in G | \omega (x) > \lambda \}.
\]

\begin{enumerate}
\def\labelenumi{(\roman{enumi})}
\setcounter{enumi}{1}
\tightlist
\item
  is obvious. We prove (i). Clearly \(G(a_{\lambda}) = \bigcap_{\mu < \lambda} G(a_{\mu})\) and \(G = \bigcup_{\lambda > 0} G(a_{\lambda})\) . On the other hand, if \(x \in G(a_{\lambda})\) and \(y \in G(a_{\mu})\) , then
\end{enumerate}

\[
x ^ {[ - 1 ]}. y ^ {[ - 1 ]}. x. y \in G (a _ {\lambda + \mu})
\]

by Proposition 6 applied with \(a = a_{\lambda}, b = a_{\mu}, c = c' = a_{\lambda + \mu}\) .

By Chapter II, § 4, no. 4, we can form the group gr(G) associated with the group G, with the central filtration \((G(a_\lambda))\). Writing \(G_\lambda = G(a_\lambda)/G(a_\lambda^+)\) for all \(\lambda > 0\), we obtain \(gr(G) = \bigoplus_{\lambda>0} G_\lambda\). Recall (loc. cit., Proposition 1) that the commutator in G allows us to define a bracket in gr(G) with which gr(G) is a Lie algebra, as follows: if \(x \in G_\lambda\) and \(y \in G_\mu\), choose a representative x of \(\bar{x}\) in \(G(a_\lambda)\) and a representative y of \(\bar{y}\) in \(G(a_\mu)\); then \([\bar{x},\bar{y}]\) is the class of \(x^{[-1]}.y^{[-1]}.x.y \in G(a_{\lambda+\mu})\) in \(G_{\lambda+\mu}\). By Proposition 6, applied with \(a = a_\lambda, b = a_\mu, c = a_{\lambda+\mu}, c' = a_{\lambda+\mu}\), we see that \([\bar{x},\bar{y}]\) is also the class of \([x,y]\) in \(G_{\lambda+\mu}\). Thus, when G is considered as a Lie subalgebra of L(G) = \(K^r\), filtered by the \(G(a_\lambda)\), the associated graded Lie algebra (Chapter II, § 4, no. 3) is equal to gr(G).

